\section{Related Work}
\label{sec:related}

\para{Adversarial attacks on LLMs.}
\citet{zou2023universal} introduced the Greedy Coordinate Gradient (\gcg) attack, which optimizes adversarial suffixes in discrete token space to bypass safety alignment. These suffixes transfer across models, achieving high attack success rates even on commercial APIs. \citet{zhu2023autodan} proposed \autodan, which uses a hierarchical genetic algorithm to produce more \emph{readable} adversarial prompts, bridging the gap between \gcg gibberish and human-crafted jailbreaks. \citet{wei2023jailbroken} taxonomized jailbreak failure modes into two categories---competing objectives and mismatched generalization---providing a theoretical framework for understanding why safety training fails on out-of-distribution inputs. Unlike these works, which focus on attack \emph{success}, we study whether models can \emph{explain} the attacks that manipulate them.

\para{LLM comprehension of nonsense.}
\citet{cherepanova2024talking} systematically studied \babel---gibberish prompts crafted via \gcg that compel specific coherent outputs. They found that harmful text is \emph{easier} to elicit than benign text (81\% vs.\ 35--66\% success on Vicuna-7B), that Babel prompts cluster separately from random tokens in representation space despite similar perplexity, and that removing a single token breaks over 70\% of prompts. \citet{erziev2025sens} discovered that LLMs can understand English instructions encoded in visually incomprehensible Unicode sequences, with Claude models achieving $\sim$30\% understanding rates. Both papers study model \emph{responses} to nonsense; we extend this line of work by evaluating model \emph{self-explanation}---whether models can articulate \emph{what} they understood and \emph{why} they responded.

\para{Perplexity-based defenses.}
\citet{alon2023detecting} proposed using \gpttwo perplexity to detect \gcg-style adversarial suffixes, showing that these suffixes have exceedingly high perplexity compared to natural text. \citet{jain2023baseline} evaluated baseline defenses including perplexity filtering, paraphrasing, and retokenization, finding that windowed perplexity blocks 80\% of white-box attacks with $\sim$7\% false positives. \citet{robey2023smoothllm} introduced \smoothllm, which exploits the fragility of adversarial suffixes through randomized character perturbation. These defenses treat perplexity as a binary signal (adversarial vs.\ natural). We instead study how model comprehension \emph{varies continuously} across the perplexity spectrum, revealing that the relationship between perplexity and explanation quality is non-monotonic.

\para{Mechanistic interpretability of adversarial prompts.}
Recent work has begun to explain \emph{how} adversarial suffixes manipulate model behavior at a mechanistic level. \citet{arditi2024refusal} showed that refusal behavior in LLMs is mediated by a single direction in activation space, suggesting that adversarial suffixes may work by suppressing this direction. \citet{chen2025attention} demonstrated that universal jailbreak suffixes function as attention hijackers, redirecting attention patterns away from safety-relevant tokens. \citet{wang2024noise} showed that adversarial suffix embeddings can be translated back into coherent text, suggesting that gibberish encodes interpretable information at the embedding level. Our work complements these mechanistic analyses with a behavioral study: rather than probing internal representations, we test whether models can articulate their own (mis)understanding.
